% Lösungen Zone (zusammenbau v0.1)
\einheitenkopf[][Kennst du schon]{Das kennst du schon}

Zuordnung: 1 → das Lösungswerkzeug aller Einheiten · 2 → die Geradenzeilen in Einheit eins und die Tangentenbedingungen in Einheit zwei · 3 → die Ansatzwahl in Einheit eins · 4 → die Ansatzverkürzung in Einheit eins · 5 → die Steigungsbedingungen in Einheit zwei · 6 → die Ansätze in Einheit drei · 7 → das Lösungswerkzeug aller Einheiten · 8 → das Lösungswerkzeug aller Einheiten

\erg{1}{a) $x = 4$, $y = 3$ \quad b) $a + b = -1$ und $4a + 2b = 6$, also $a = 4$, $b = -5$}
\erg{2}{a) $m = 2$, $n = 3$, also $y = 2x + 3$ \quad b) $m = \frac{2 - 5}{4 - (-2)} = -0{,}5$; $n = 4$; $y = -0{,}5x + 4$}
\erg{3}{a) $S(4|2)$ \quad b) $S(-2|3)$, nach unten geöffnet}
\erg{4}{a) achsensymmetrisch zur $y$-Achse – nur gerade Exponenten, die Zahl zählt als $x^0$. \quad b) punktsymmetrisch zum Ursprung – $x^5$, $x^3$ und $x$ haben ungerade Exponenten.}
\erg{5}{a) $f'(x) = 12x^2 - 2$ \quad b) $f'(x) = x^2 - 3x$; $f'(-1) = 4$}
\erg{6}{a) Amplitude $5$ \quad b) Periode und Amplitude: $4$ und $2$}
\erg{7}{a) Beim Subtrahieren wurde $-(-b)$ nicht zu $+b$; $b$ fällt weg. Richtig: $3a = 6$, also $a = 2$, $b = 1$.}
\erg{8}{a) Subtrahieren: $3a = 9$, also $a = 3$, $b = 3$}
